\chapter{1980年全国卷（文）}

\begin{problem}
化简： \(\frac{1 - 3\i}{3 - 2\i}\).

\end{problem}

\begin{problem}
解方程组： \(\begin{cases}2x - 3y - z = 5\\ 4x + 2y + 3z = 9\\ 3x + 2y = -1\end{cases}\).

\end{problem}

\begin{problem}
用解析法证明：直径所对的圆周角是直角.

\end{problem}

\begin{problem}
某地区 1979 年的轻工业产值占工业总产值的 20\%，要使 1980 年的工业总产值比上一年增长 10\%，且使 1980 年的轻工业产值占工业总产值的 24\%，问 1980 年轻工业产值应比上一年增长百分之几？

\end{problem}

\begin{problem}
设 \(\frac{3\pi}{4} < \theta < \frac{5\pi}{4}\)，化简： \(\frac{\sqrt{\cos\frac{\pi}{4}\sin\left(\frac{3\pi}{4}-\theta\right)}\left[\sin(\pi-\theta)-\sin\left(\theta-\frac{\pi}{2}\right)\right]}{\sin\left(\theta+\frac{\pi}{4}\right)}\).

\end{problem}

\begin{problem}
\begin{enumerate}
\item 若四边形 \(ABCD\) 的对角线 \(AC\) 将四边形分成面积相等的两个三角形，证明：直线 \(AC\) 必平分对角线 \(BD\)；
\item 写出第1小题的逆命题，这个逆命题是否正确？为什么？
\end{enumerate}

\end{problem}

\begin{problem}
如图，长方形框架 \(ABCD - A'B'C'D'\) 三边 \(AB\)、\(AD\)、\(AA'\) 的长分别为 6、8、3.6，\(AE\) 与底面的对角线 \(B'D'\) 垂直于 \(E\).

\begin{center}
\includegraphics[page=1,trim=34.93cm 21.93cm 1.57cm 4.49cm,clip,width=6.0cm]{source/普通高考/1980/1980文.pdf}
\end{center}

\begin{enumerate}
\item 证明 \(A'E \perp B'D'\)；
\item 求 \(AE\) 的长.
\end{enumerate}

\end{problem}

\begin{problem}
\begin{enumerate}
\item 把参数方程 \(\begin{cases}x = \sec t\\ y = 2\tan t\end{cases}\) （\(t\) 为参数）化为直角坐标方程，并画出方程的曲线的略图；
\item 当 \(0\leqslant t < \frac{\pi}{2}\) 及 \(\pi \leqslant t < \frac{3\pi}{2}\) 时，各得到曲线的哪一部分？
\end{enumerate}

\end{problem}
